\chapter{Annexe B — Sources de données}
\label{ann:sources}

Chaque fait publié par le service provient d'une source identifiée. Ce tableau
est la version humaine de ce que le contrat porte champ par champ.

\begin{table}[h]
\centering
\small
\begin{tabularx}{\textwidth}{@{}p{2.9cm}p{3.3cm}p{2.6cm}X@{}}
\toprule
\textbf{Donnée} & \textbf{Producteur} & \textbf{Résolution} & \textbf{Usage dans le service} \\
\midrule
Aléas réglementaires & Géorisques (BRGM / MTE) & point dans polygone & présence et preuve de chaque aléa \\
Emprises de crue TRI & Géorisques (WFS) & point dans polygone & classe de crue et niveau d'eau du scénario \\
Fiche bâtiment & BDNB (CSTB) & bâtiment & 139 champs verbatim : matières, année, hauteur, empreinte \\
Adresse & BAN / Géoplateforme (IGN) & bâtiment (interpolation) & géocodage, point de jointure \\
Hydrographie & BD TOPO (IGN) / Sandre & tronçon & cours d'eau, distance, trajet de l'eau \\
Altitude & Altimétrie (IGN) & point & altitude au sol du bâtiment \\
Sol & Aléa argile (BRGM) & \textbf{commune} & sensibilité au retrait-gonflement \\
Météo prévue & Open-Meteo & point & pluie horaire, rafales, humidité \\
Vigilance crues & Vigicrues & tronçon & alerte en cours à proximité \\
Photo de terrain & Panoramax (IGN / DINUM) & prise de vue & contexte visuel et crédit \\
Prose du rapport & Fournisseur configurable & — & reformulation sous validation \\
\bottomrule
\end{tabularx}
\caption{Sources, producteurs et finesse réelle de chaque donnée.}
\end{table}

\section{Les règles de provenance}

\begin{itemize}
  \item Tout objet d'aléa porte \code{source}, \code{resolution} et
    l'attribution de licence ouverte (LO 2.0) du producteur.
  \item L'attribution est accompagnée d'un millésime : une donnée sans date de
    mise à jour n'est pas exploitable dans un contexte de conformité.
  \item Une source qui ne répond pas laisse une trace dans le contrat (absence
    déclarée) ; elle ne disparaît pas silencieusement.
  \item L'aléa argile est \textbf{toujours} publié en estimation communale tant
    qu'aucune source vectorielle n'est câblée. Le changer exige un amendement de
    la constitution du projet.
\end{itemize}

\begin{invariant}[title={Deux seuils, jamais confondus}]
Un seuil de 0,3 m gouverne les \emph{dommages} (compteurs d'objets touchés) ;
la carte, elle, représente l'eau dès la première profondeur positive. Confondre
les deux ferait disparaître de la carte des zones réellement inondées.
\end{invariant}

\section{Vérification de l'installation}

\begin{lstlisting}[language=bash,caption={Démarrage local et contrôle de vie.}]
# Backend
cd backend && ./venv/Scripts/python.exe -m uvicorn app.main:app --port 8000
# Frontend
cd frontend && npm run dev            # http://localhost:5173

# Controle de vie
curl http://127.0.0.1:8000/health     # {"status":"ok"}

# Suite de tests hors-ligne
cd backend && ./venv/Scripts/python.exe -m pytest tests/ -q

# Compilation de ce rapport
tectonic -X compile rapport/main.tex
\end{lstlisting}
